\documentclass[11pt]{article}
\usepackage[utf8]{inputenc}
\usepackage{amsmath,amssymb}
\usepackage[margin=1in]{geometry}
\usepackage{hyperref}
\emergencystretch=3em

\title{Symmetric normal boxes: the bare parity identity}
\author{Claude, with Daniel Murfet}
\date{2026-09-07}

\begin{document}
\maketitle
\sloppy

\begin{abstract}
Interior normal coordinates range over $(-1,1]$ rather than $(0,1]$. For the bare monomial integral
the coordinate reflections give the exact identity
\[
\int_{(-1,1]^d}u^he^{-\beta Nu^{2k}}\,du=\prod_i\big(1+(-1)^{h_i}\big)\int_{(0,1]^d}u^he^{-\beta Nu^{2k}}\,du
\]
(\texttt{monomialSymReal\_eq}): each odd exponent kills the integral, each even one doubles it ---
the paper's parity factor $(1+(-1)^{h_i})$ without the $\tfrac12$ of an averaged measure. Hence the
symmetric-box integral has the mixed-ratio asymptotic with the constant multiplied by the parity
product (\texttt{monomialSymReal\_tendsto}). Zero \texttt{sorry}/\texttt{axiom}.
\end{abstract}

\section{The things formalised}
{\small
\begin{verbatim}
theorem integrableOn_pi_box_marginal (d S F) (hF : IntegrableOn F (piBox (d + 1) S)) :
    IntegrableOn (fun a => ∫ b in piBox d S, F (Fin.cons a b)) S
theorem integral_Ioc_symm_of_parity (g) (h : ℕ) (hg : IntegrableOn g (Ioc (-1) 1))
    (hsym : ∀ a, g (-a) = (-1) ^ h * g a) :
    ∫ a in Ioc (-1) 1, g a = (1 + (-1) ^ h) * ∫ a in Ioc 0 1, g a
def symBox (d : ℕ) : Set (Fin d → ℝ) := piBox d (Ioc (-1) 1)
def monomialSymReal (d) (h k) (β N) : ℝ :=
  ∫ x in symBox d, (∏ i, x i ^ h i) * exp (-(β * N * ∏ i, x i ^ (2 * k i)))
theorem monomialSymReal_eq (β) : ∀ d h k N,
    monomialSymReal d h k β N = (∏ i, (1 + (-1 : ℝ) ^ h i)) * monomialBoxReal d h k β N
theorem monomialSymReal_tendsto ... : Tendsto (fun N => monomialSymReal (d + 1) h k β N
      / (N ^ (-l) * log N ^ (multCount (ratioExp h k) l - 1))) atTop
      (𝓝 ((∏ i, (1 + (-1 : ℝ) ^ h i)) * monomialMixedConst h k l β))
\end{verbatim}
}

\section{Mathematics}
Induction on the dimension: the Bochner peel (unit~188) exposes the first coordinate on both boxes,
the inner symmetric integral is the $d$-dimensional one at parameter $Na^{2k_0}$ (induction
hypothesis), and the one-dimensional reflection identity handles the peeled coordinate, whose
integrand $a\mapsto a^{h_0}G(a^{2k_0})$ has parity $(-1)^{h_0}$. Integrability of the peeled marginal
comes from \texttt{Integrable.integral\_prod\_left} along the same measure-preserving transport.

\section{Paper-facing meaning}
This is the bare version of the paper's interior-coordinate reduction (parity factors
$\frac{1+(-1)^{h_i+\gamma_i}}2$ up to the normalisation of the measure). For a general amplitude the
reflections produce the signed combination $\sum_\sigma\prod_i\sigma_i^{h_i}\eta(\sigma\cdot u)$, not a
scalar factor; an odd exponent then cancels only the corner coefficient, not the whole integral.
That reduction is a natural follow-up.

\section{Lean notes}
\begin{itemize}
\item The reflection identity is cleanest as a statement about an integrable $g$ with pointwise parity,
proved with \texttt{integral\_add\_adjacent\_intervals} and \texttt{integral\_comp\_neg} on interval
integrals, converting to \texttt{Ioc} integrals with \texttt{integral\_of\_le} at the ends.
\item \texttt{IntegrableOn.congr\_fun} with the peel identity transfers marginal integrability to the
rewritten integrand; its pointwise goal is a beta-redex (\texttt{beta\_reduce}).
\item \texttt{unitBox d = piBox d (Ioc 0 1)} and \texttt{symBox d = piBox d (Ioc (-1) 1)} are
\texttt{rfl}; state them as \texttt{have} equations to \texttt{rw} into the peel lemma's shape.
\end{itemize}

\end{document}
